\section{Setup}
\label{sec:setup}

\subsection{One loss for all contrastive methods}
A batch holds $B$ images. Each is augmented twice, which gives $2B$ views with $\ell_2$-normalised
projections $z_1,\dots,z_{2B}$. Let $n(i)$ be the image from which view $i$ was made. Every image $n$
carries a \emph{group} $g_n$, and two different views are positives when their images are in the same
group:
\begin{equation}
\label{eq:loss}
\mathcal{L} \;=\; \frac{1}{2B}\sum_{i=1}^{2B} \frac{-1}{|P(i)|}\sum_{p\in P(i)}
\log\frac{\exp(z_i^\top z_p/\tau_i)}{\sum_{a\neq i}\exp(z_i^\top z_a/\tau_i)},
\qquad P(i)=\{\,p\neq i : g_{n(p)} = g_{n(i)}\,\}.
\end{equation}
The temperature $\tau_i\in\{\tau_\ell,\tau_u\}$ depends on the anchor: $\tau_\ell$ if image $n(i)$ is used
with its label and $\tau_u$ otherwise. Three choices of groups give the three methods.
\begin{itemize}[leftmargin=1.4em]
\item \textbf{\simclr} \citep{chen2020simclr}: no label is used and every image is its own group
($g_n=n$). $P(i)$ is the other view of the same image, and \cref{eq:loss} is the NT-Xent loss with
temperature $\tau_u$.
\item \textbf{\supcon} \citep{khosla2020supcon}: every image in the batch is labeled and its group is its
class ($g_n=y_n$). \Cref{eq:loss} is the $\mathcal{L}^{\mathrm{sup}}_{\mathrm{out}}$ loss of
\citet{khosla2020supcon} with temperature $\tau_\ell$.
\item \textbf{\union} (\supcon with unlabeled images): labeled images use their class and unlabeled images
are their own group. An unlabeled image is still contrasted against every other view in the batch, so it
supplies negatives and its own instance-level positive; it never claims a class.
\end{itemize}
We use $\tau_\ell=0.1$ and $\tau_u=0.5$, the standard CIFAR values of \supcon and \simclr
\citep{khosla2020supcon}. Because the temperature depends on the anchor, the loss of \union is not symmetric
in a labeled--unlabeled pair: the pair is weighted with $\tau_\ell$ when the labeled view is the anchor and
with $\tau_u$ when the unlabeled one is. \union is close to SuNCEt \citep{assran2020suncet} and to the
generalised loss of \citet{inoue2020semi}. Its temperatures and the share of labeled images per batch (25\%)
were chosen on a validation split of the training set before any test number was computed
(\cref{app:pilot}). Temperature is thus a second difference between the label-free and the label-using
rows of our design; \cref{sec:plateau} shows that swapping the temperatures does not change the picture.
Our implementation of \cref{eq:loss} agrees with the official \supcon code to $10^{-6}$ (\cref{app:impl}).

\subsection{The \texorpdfstring{$2\times2$}{2x2} design}
With a labeled fraction $f$ (10\% unless stated), the label-scarce setting couples two things: the
\emph{labels} that the loss may use and the \emph{images} in the pretraining pool. \Cref{tab:design} varies
them independently. Comparing the rows gives the effect of the labels at a fixed pool. Comparing the columns
gives the effect of the unlabeled images at a fixed use of labels. The labeled subsets are class-balanced
and nested across fractions for a given seed.

\begin{table}[t]
\centering
\small
\caption{The $2\times2$ design. Each cell is one pretraining method; all four receive the same number of
gradient steps.}
\label{tab:design}
\begin{tabular}{lcc}
\toprule
 & pool = the labeled images only & pool = all images \\
\midrule
labels not used in the loss & \simclrsub & \simclr \\
labels used in the loss & \supcon & \union \\
\bottomrule
\end{tabular}
\end{table}

\subsection{Compute accounting}
\label{sec:compute}
\paragraph{Steps, not epochs}
Every contrastive run performs $T=E\cdot\lfloor N/B\rfloor$ gradient steps with batch size $B=512$, where $N$
is the size of the \emph{full} training set, whatever the size of the method's own pool. On CIFAR
($N=50{,}000$) this is $97E$ steps of 1,024 augmented views. We report $E$ from 10 to 1,000, that is, 970 to
97,000 steps, and always state budgets in steps. A method whose pool has 5,000 images simply cycles through
it ten times as often.

The distinction matters. If each method were instead given ``100 epochs'' over its own pool, \supcon on
5,000 images would run $100\cdot\lfloor 5{,}000/512\rfloor = 900$ steps against 9,700 for \simclr on all
images, and the effect of supervision would be mixed with an eleven-fold difference in compute. In our
experiments this alone moves \supcon by \R{d:steps900}\pt (\R{supcon_lab_f0.1_steps900:lp50} at 900 steps
against \R{supcon_lab_f0.1_e100:lp50} at 9,700). We do not claim that the confusion is widespread: in the
audit of \cref{app:audit}, 7 of 17 applicable papers matched or accounted for the difference, 2 compared
arms at equal epochs over pools of different size, and 5 did not state the budget of one arm in a way a
reader can check. Stating budgets in steps or in images processed removes the ambiguity.

\paragraph{Cost per step}
\Cref{tab:compute} lists, for every method, the number of augmented views, forward passes and backward
passes per step, the number of steps at the reference budget and the wall-clock time on one NVIDIA A40.
The three contrastive methods have identical cost. BYOL adds a forward pass through its target network. The
semi-supervised baselines are given the same number of processed views as a 9,700-step contrastive run
unless stated; their per-step composition differs, so we report it explicitly.

\begin{table}[t]
\centering
\small
\caption{Compute accounting on CIFAR-10 with a ResNet-18. ``Reference budget'' is the budget labeled 9,700
steps in the text. FixMatch processes 64 labeled, 448 weakly and 448 strongly augmented unlabeled views per
step; MixMatch needs an extra forward pass to guess labels. Minutes are wall-clock training time on one
NVIDIA A40 with no other job on the GPU.}
\label{tab:compute}
\begin{adjustbox}{max width=\linewidth}
% generated by scripts/make_assets.py -- do not edit
\begin{tabular}{lcccrcc}
\toprule
 & \multicolumn{3}{c}{per step} & & \multicolumn{2}{c}{reference budget} \\
\cmidrule(lr){2-4}\cmidrule(lr){6-7}
method & augmented views & forward passes & backward passes & steps & forward passes ($10^6$) & minutes \\
\midrule
\simclr, \supcon, \union & 1,024 & 1,024 & 1,024 & 9,700 & 9.9 & 12 \\
BYOL & 1,024 & 2,048 & 1,024 & 9,700 & 19.9 & 22 \\
SimSiam & 1,024 & 1,024 & 1,024 & 9,700 & 9.9 & 15 \\
cross-entropy & 512 & 512 & 512 & 9,700 & 5.0 & 6 \\
FixMatch & 64 + 448 + 448 & 960 & 960 & 10,347 & 9.9 & 24 \\
MixMatch & 64 + 2$\times$64 & 320 & 192 & 51,733 & 16.6 & 37 \\
fine-tuning on the 5k (FT-5k) & 128 & 128 & 128 & 2,344 & 0.3 & 1 \\
\bottomrule
\end{tabular}

\end{adjustbox}
\end{table}

\subsection{Evaluation protocols}
\label{sec:protocols}
The encoder is frozen (except in fine-tuning) and its average-pooled features are used in four ways. We
name the protocols once and use the names throughout.
\begin{itemize}[leftmargin=1.4em]
\item \textbf{\lpall, the oracle-label probe.} A linear classifier trained on the features of all
50,000 training images with all their labels \citep{chen2020simclr}. It is the standard measure of
representation quality, but it uses ten times more labels than the setting provides. We also report a
weighted $k$NN classifier with the same 50,000 labels ($k$NN-50k; \citealp{wu2018instance}).
\item \textbf{\lpfive.} The same linear probe trained only on the 5,000 labeled images, the ones the
label-using methods saw during pretraining.
\item \textbf{\ftfive.} Fine-tuning of the whole network on the 5,000 labeled images. The default recipe
starts from the \lpfive solution \citep{kumar2022lpft} and trains for 60 passes over the labeled images
with learning rate 0.01. \Cref{sec:protocol} replaces it by a recipe selected on a split of the 5,000
labels.
\item \textbf{FixMatch-init.} FixMatch \citep{sohn2020fixmatch} started from the pretrained encoder, using
the 5,000 labels and all images without labels.
\end{itemize}
\lpfive, \ftfive and FixMatch-init respect the label budget, and we refer to them together as
\emph{label-budget protocols}. For other fractions we write LP-all and LP-$f$. On STL-10 only 5,000 training
labels exist, so there is no oracle-label probe and every protocol is a label-budget protocol. We consider
the pair (oracle-label probe, label-budget protocol) to be part of the method of this paper: the first says
what an encoder has learned, the second what a practitioner can extract from it, and \cref{sec:protocol}
shows that the two can disagree about which encoder is better.

The linear probes are $\ell_2$-regularised multinomial logistic regressions on standardised features, fitted
with L-BFGS; the regularisation strength is selected on a stratified 20\% of the labels available to the
probe. No design decision in this paper was taken on test data.

\subsection{Data, models and training}
\paragraph{Datasets}
CIFAR-10 and CIFAR-100 \citep{krizhevsky2009cifar} have 50,000 training and 10,000 test images of 32 px; we
use $f=10\%$ (5,000 labels) unless stated and fractions from 1\% to 100\% in \cref{sec:fractions}. STL-10
\citep{coates2011analysis} has 5,000 labeled training images, 8,000 test images and 100,000 unlabeled images
at 96 px. Its unlabeled set comes from a broader distribution than the ten classes, so it tests the
``all images'' column with data that are not curated to match the labels.

\paragraph{Models and optimisation}
The encoder is a ResNet-18 \citep{he2016resnet} with a $3\times3$ stride-1 first convolution and no max-pooling
for 32-px images; for STL-10 the max-pooling layer is kept. A ResNet-50 is used in \cref{sec:datasets}. The
projection head is linear--batch-norm--ReLU--linear with 128 outputs. We train with SGD (momentum 0.9,
weight decay $10^{-4}$, learning rate 0.25 for batch 512, 5\% linear warm-up, cosine decay), the recipe of the
official \supcon code, in mixed precision with the loss computed in float32. The augmentation is the CIFAR
recipe of \citet{chen2020simclr}: random resized crop (scale 0.2--1), horizontal flip, colour jitter of
strength 0.5 applied with probability 0.8, and random grayscale with probability 0.2. Baselines (BYOL,
SimSiam, cross-entropy, FixMatch, MixMatch) and all remaining hyper-parameters are described in
\cref{app:hparams}.

\subsection{Statistics}
\label{sec:stats}
The main configurations are run with five seeds at 9,700 and 19,400 steps and with three seeds elsewhere;
the seed sets the labeled subset, the initialisation and the augmentation stream. We report the mean and
the standard deviation over seeds. Differences between two configurations are tested with a
\emph{hierarchical bootstrap} \citep{efron1993bootstrap}. One replicate first resamples the seeds with
replacement (jointly for the two configurations when they share seeds, because runs with the same seed share
the labeled subset) and then resamples the test images with replacement, using the same resampled images
for both configurations. We draw 100,000 replicates of the difference in mean accuracy and report its
percentile 95\% interval and a two-sided $p$-value. The \R{stats:ntests} planned comparisons of
\cref{tab:stats} form one family, and their $p$-values are adjusted with the Holm--Bonferroni procedure
\citep{holm1979simple}. A statement that two methods perform the same is backed by an equivalence test: two
one-sided tests \citep{schuirmann1987comparison} on the bootstrap distribution with a margin of $\pm1$ point.
With three to five seeds the between-seed variance is estimated coarsely, so we treat intervals as
approximate and do not interpret differences below about one point. Results from a single seed are marked
as such.
